% !TEX root = userguide.tex

\def\headdate{27th July 2023}

\section{Arbitrary Lagrangian Eulerian formulation}
\label{sec:ale}

If we have moving boundaries and/or interfaces, the mesh will change as a
function of time. If the mesh moves independently from the material, we
use the so-called Arbitrary Lagrangian Eulerian (ALE)
formulation. In \xfem\ the mesh does not move and
we will use a special ALE version called tALE.
In this section examples using the ALE
formulation will be given after the theory in the next section.

\subsection{Basics}

In order to describe inhomogeneous functions we need coordinates as function
argument. For example, consider a quantity $f$ specified by a function:
\begin{equation}
 f=\hat f(\vek x,t)
\end{equation}
where the arguments are the
(fixed) spatial coordinates (position vector) $\vek x$ and time $t$. This is
called the spatial or Eulerian formulation. However, it is also possible to use
a moving curvilinear coordinate system given by
\begin{equation}
   \vek x = \hat{\vek x}(\col\xi,t)
    \label{eq:mappingxi}
\end{equation}
where $\col\xi=(\xi_1,\xi_2,\xi_3)$ are the curvilinear coordinates and
$\hat{\vek x}$ is the function that maps the coordinates $\col\xi$ onto the
spatial coordinates $\vek x$. Note, that the mapping is a function of time
(moving coordinates). We now can find a new function
with $\col\xi$ as an argument to define the quantity $f$:
\begin{equation}
   f=\hat f(\vek x,t)=\hat f(\hat{\vek x}(\col\xi,t),t)\equiv \bar
f(\col\xi,t)
  \label{eq:fexpr}
\end{equation}
In the following we will call $\col\xi$ the grid coordinates, or simply the
grid.

We have the following extreme cases:
\begin{itemize}
   \item The mapping function $\hat{\vek x}$ is independent of time. The grid is
``fixed'' and basically we have a Eulerian formulation, where the grid is just
a curvilinear coordinate system to describe space.
   \item The grid coordinates $\col\xi$ move with the material. Thus each
material point can be identified with the unique coordinates $\col\xi$
(``material label''). This is called a Lagrangian formulation.
\end{itemize}
If the grid can move independently it is called an Arbitrary Lagrangian
Eulerian (ALE) formulation. The ALE formulation helps to minimize mesh
distortion while still
being able to track material boundaries similar to a Lagrangian
formulation.

We can define the grid time derivative:
\begin{equation}
   \frac{\delta f}{\delta t} = \pderiv{f}{t}\Big|_{\col\xi=\text{constant}}
                             = \pderiv{\bar f}{t}
\end{equation}
Hence, the grid time derivative describes the change in time of the quantity
$f$ as seen from a fixed grid point $\col\xi$.	In order to find an expression
for the grid time derivative we differentiate Eq.~\eqref{eq:fexpr} in time
using the chain rule:
\begin{equation}
   \frac{\delta f}{\delta t} = \pderiv{\hat f}{t}+\pderiv{\hat{\vek
x}}{t}\cdot\nabla \hat f
\end{equation}
Here, the first term on the right-hand side is the local time derivative
(usually denoted by $\plderiv{f}{t}$), the first factor in the second term is
the grid velocity (denoted by $\vek u_\text{g}$) and the second factor is the
spatial gradient of $f$. Therefore, we get
\begin{equation}
   \frac{\delta f}{\delta t} = \pderiv{f}{t}+\vek u_\text{g}\cdot\nabla f
  \label{eq:gridderiv}
\end{equation}
For the special cases we get:
\begin{itemize}
   \item Eulerian formulation, $\vek u_\text{g}=\vek 0$:
\begin{equation}
    \frac{\delta f}{\delta t} = \pderiv{f}{t}
\end{equation}
   \item Lagrangian formulation, $\vek u_\text{g}=\vek u$ (velocity of the
fluid):
\begin{equation}
    \frac{\delta f}{\delta t} = \pderiv{f}{t}+\vek u\cdot\nabla f\equiv\dot f
  \label{eq:substan}
\end{equation}
where $\dot f$ is the material (or substantial) derivative.
\end{itemize}
Combining Eqs.~\eqref{eq:gridderiv} and \eqref{eq:substan}$_2$ gives
\begin{equation}
    \dot f=\frac{\delta f}{\delta t}+ (\vek u-\vek u_\text{g})\cdot\nabla f
  \label{eq:substan2}
\end{equation}
which is quite useful for transforming the physical equations into a suitable
form for implementation.

In conventional ALE the grid will be identified with the moving mesh and
the grid
derivative can be described by the time derivative of the nodal values as in the
Eulerian formulation. In the convection terms the velocity needs to be
changed (subtraction of the grid velocity). 

In \xfem\ (see Sec.~\ref{sec:xfem}) the mesh is usually not moving in time, but
the boundaries/interfaces often are. Here we use also Eq.~\eqref{eq:substan2}
to handle material derivatives. Since the mesh is not moving, we cannot
identify the grid with the computational mesh. Therefore we will use a so-called
temporary ALE (tALE) mesh for that. A tALE mesh describes a mapping between two
consecutive time steps. For a first-order time discretization we only need 
one tALE mesh, which is renewed every time step. For a second-order multi-step
discretization we need two tALE meshes, with one renewal every time step.
The grid velocity $\vek u_\text{g}$ is defined by (backwards) differentiation
of the grid position.
Note, that the tALE meshes are different from the computational mesh.
The tALE technique has been developed by
Choi \cite{Choi11b}. See also \cite{Choi11c,choi_alignment_2012} for a detailed description
of the tALE technique as implemented in \tfem.

\subsection{Oscillating planar Couette flow with a freely suspended particle
in a viscoelastic liquid}
\label{sec:couette_ale}

(This example has been developed by Gaetano d'Avino of the University of
Naples)

We consider a planar Couette problem (shear flow between two moving walls) with
a (off-center) freely suspended particle. 
The curves describing the domain are plotted in Fig.~\ref{fig:curves_ale1}.
\begin{figure}[htbp!]
   \begin{center}
   \includegraphics[scale=0.5]{ale1_curves}
   \end{center}
   \caption{Points and curves of the oscillating
     Couette flow problem with a suspended particle.}
    \label{fig:curves_ale1}
\end{figure}
The walls (at curves \texttt{C1} and \texttt{C3}) 
are moving in an oscillatory way and 180~degrees out of phase. 
In horizontal direction periodical boundary conditions are assumed.
In the lower left corner the pressure value is set to zero (Dirichlet).
The viscoelastic model is a Giesekus fluid.

The initial mesh is generated using Gmsh (see Sec.~\ref{sec:readgmshmesh})
and is shown in Fig.~\ref{fig:mesh_ale1}.
\begin{figure}[htbp!]
   \begin{center}
   \includegraphics[scale=0.5]{ale1_mesh}
   \end{center}
   \caption{Mesh of the oscillating
     Couette flow problem with a suspended particle.}
    \label{fig:mesh_ale1}
\end{figure}

The example file is \texttt{ale1.f90}\footnote{The example does not run
currently. Combining the ALE approach with time-integration schemes requiring
multiple evaluations of the convection term at different times in a single time
step does not take the proper mesh coordinates into account. Use example \texttt{ale2.f90} instead.}
and can be obtained by typing at the
command line
\medskip
\begin{quote}
   \texttt{getexample ale}
\end{quote}

Some notes on the source:
\begin{itemize}
   \item
     The periodical boundary conditions are imposed using
     constraints (see Chap.~\ref{ch:constraints}).
   \item 
     The explicit stress formulation with DEVSS-G/SUPG is used (see
     Sec.~\ref{sec:explicitstress}). 
   \item The position of the particle is obtained using an explicit
     second-order Adams-Bashforth scheme.
   \item The grid velocity is defined by solving a Laplace equation
   \begin{equation}
      \nabla\cdot(k^e\nabla\vek u_\text{g})=\vek 0 \label{eq:laplace_ale}
   \end{equation}
     where $\vek u_\text{g}$ is equal to the particle velocity on the surface
     of the particle (curve \texttt{C5}). On all other boundaries $\vek
     u_\text{g}$ is set to zero.
    In Eq.~\eqref{eq:laplace_ale} 
   $k^e$ is a function that controls the deformation of the domain.
   As suggested in \cite{Hu2001} we take $k^e$ 
   constant on an element and equal to the inverse of the
     area/volume of the element. In this way the deformation of the mesh is
largely performed by the big elements and the small ones around the solid
particle just move almost rigidly with the particle.
   \item
    The conformation tensor equation (constitutive model)
    is written in the ALE grid using
    Eq.~\eqref{eq:substan2}.
    A Crank-Nicolson/Adams-Bashforth (second-order)
    with velocity prediction time discretization
    (see Eq.~\eqref{eq:explicitCrankAB2}) is used to discretize the grid time
    derivative.
    \item \texttt{vtk} files are written to make a movie with Paraview.
\end{itemize}

After running the example
the results in Fig.~\ref{fig:ale1_position} for the vertical position of the
particle is obtained.
\begin{figure}[htbp!]
   \begin{center}
   \includegraphics[scale=0.8]{ale1_position}
   \end{center}
   \caption{Vertical position of the particle as a function of time.}
    \label{fig:ale1_position}
\end{figure}
We see a migration of the particle towards the wall.

Note, that a similar problem has been described in \cite{DAvino10a}. An
important difference, however, is that in \cite{DAvino10a} the horizontal
grid velocity is taken equal to the particle velocity for \emph{all nodes}. In
this way the grid is only deformed if the particle moves vertically. This 
minimizes grid distortion, but is not applicable if multiple particles are
present.
 
A second example \texttt{ale2.f90} is available with the following differences:
\begin{itemize}
   \item 
     The implicit stress formulation with DEVSS-G/SUPG is used (see
     Sec.~\ref{sec:implicitstress}). 
   \item
    A second-order semi-implicit Gear scheme with conformation prediction
    (see Eq.~\eqref{eq:Gearimp}) is used to discretize the grid time
    derivative.
\end{itemize}

\subsection{Oscillating planar Couette flow with a freely suspended particle
in a Newtonian liquid. Effect of particle and fluid inertia included.}
\label{sec:couette_ale_inertia}

(This example has been developed by Gaetano d'Avino of the University of
Naples)

We again
consider a planar Couette problem (shear flow between two moving walls) with
a (off-center) freely suspended particle. The example is similar to the example
in Sec.~\ref{sec:couette_ale}, but now the fluid is Newtonian and
inertia is included for both the particle and the fluid. 

The example files are \texttt{ale3.f90} and \texttt{ale4.f90}.
and can be obtained by typing at the
command line:
\medskip
\begin{quote}
   \texttt{getexample ale}
\end{quote}

Some notes on the source:
\begin{itemize}
   \item The difference between \texttt{ale3.f90} and \texttt{ale4.f90} is the
way the particle motion is handled. In \texttt{ale3.f90} an Adams-Bashforth
scheme is used (see Figure~\ref{fig:particle_inertia_AB2}). 
In \texttt{ale4.f90}
a predictor-corrector scheme (see Figure~\ref{fig:particle_inertia_predcor}) is used.
\begin{figure}[p!]
   \begin{center}
   \includegraphics[scale=0.65]{step_inertia3}\hspace{1cm}
   \includegraphics[scale=0.65]{step_inertia4}
   \end{center}
   \caption{Time discretization of the particle motion in \texttt{ale3.f90}.
    At the initial time $t_0$ the initial position $\vek X_0$ and velocity
    $\vek U_0$ of the particle is given. In step 1 (left)
    the particle position is
    updated using an explicit Euler scheme. The velocity of the particle is
    found by solving the coupled system of equations using the new position of 
    the particle.
    In any step beyond 1 (right) a second order AB2 scheme is used to update the
    position of the particle.
    The velocity of the particle is
    found by solving the coupled system of equations using the new position of
    the particle.}
    \label{fig:particle_inertia_AB2}
\end{figure}
\begin{figure}[p!]
   \begin{center}
   \includegraphics[scale=0.65]{step_inertia1}\hspace{1cm}
   \includegraphics[scale=0.65]{step_inertia2}
   \end{center}
   \caption{Time discretization of the particle motion in \texttt{ale4.f90}.
    At the initial time $t_0$ the initial position $\vek X_0$ and velocity
    $\vek U_0$ of the particle is given. In step 1 (left)
    the particle position is
    predicted using an explicit Euler scheme. The velocity of the particle is
    found by solving the coupled system of equations using the predicted
    position of the particle.
    Using the new velocity of the particle a corrected position of the particle
    is computed with a trapezoidal rule.
    In any step beyond 1 (right) 
    the particle position is
    predicted using an extrapolation scheme. The velocity of the particle is
    found by solving the coupled system of equations using the predicted
    position of the particle.
    Using the new velocity of the particle a corrected position of the particle
    is computed with a trapezoidal rule.}
    \label{fig:particle_inertia_predcor}
\end{figure}
   The advantages of the predictor-corrector scheme are that the current
   velocity is used in the particle position and that it can be easily extended
   to problems where the updating of the interface is more complicated (see,
   for example Sec.~\ref{sec:swellNewtonian}). A disadvantage is that the
   position of the particle-fluid interface does not coincide with the mesh
   used for computing the velocities.  
   \item The grid displacement within a time step $\vek d_\text{g}$
    is defined by solving a Laplace equation
   \begin{equation}
      \nabla\cdot(k^e\nabla\vek d_\text{g})=\vek 0 \label{eq:laplace_ale2}
   \end{equation}
     $\vek d_\text{g}$ is equal to the particle displacement
     on the surface
     of the particle (curve \texttt{C5}). On all other boundaries $\vek
     d_\text{g}$ is set to zero. In Eq.~\eqref{eq:laplace_ale2} 
   $k^e$ is a function that controls the deformation of the domain.
   As suggested in \cite{Hu2001} we take $k^e$ 
   constant on an element and equal to the inverse of the
     area/volume of the element. In this way the deformation of the mesh is
largely performed by the big elements and the small ones around the solid
particle just move almost rigidly with the particle.
   \item The mesh velocity $\vek u^{n+1}_\text{g}$ at time 
    $t_{n+1}$ is defined by backwards
   differencing of the positions of the nodes of the mesh ($\vek x_\text{g}$).
   The first time step we use a first-order differentiation:
\begin{equation}
  \vek u^{n+1}_\text{g}=(\vek x^{n+1}_\text{g} - \vek x^{n}_\text{g})/\Delta t
\end{equation}
   and for steps two and further a second-order differentiation:
\begin{equation}
  \vek u^{n+1}_\text{g}=(\frac32\vek x^{n+1}_\text{g} - 2\vek x^{n}_\text{g}
            +\frac12\vek x^{n-1}_\text{g})/\Delta t
\end{equation}
  \item The inertia terms (of both fluid and particle)
        are time discretized using an implicit scheme (first
        step: Euler implicit, then second-order implicit Gear).
  \item The system is solved using a Newton-Raphson iteration scheme.
\end{itemize}

\subsection{Planar Couette flow with a freely suspended particle
in a viscoelastic liquid. Effect of particle and fluid inertia included.}
\label{sec:couette_ale_inertia_viscoelastic}

(This example has been developed by Gaetano d'Avino of the University of
Naples)

We again
consider a planar Couette problem (shear flow between two moving walls) with
a (off-center) freely suspended particle. The example is similar to the example
in Sec.~\ref{sec:couette_ale}, but now 
inertia is included for both the particle and the fluid. Also, the walls have a
constant velocity. Results on the migration of the particle in this problem
are reported in Huang et al.\ \cite{Huang97}.

The example files are \texttt{ale5.f90}, \texttt{ale6.f90}, \texttt{ale7.f90}
and \texttt{ale8.f90} and can be obtained by typing at the
command line:
\medskip
\begin{quote}
   \texttt{getexample ale}
\end{quote}

Some notes on the sources:
\begin{itemize}
   \item An Adams-Bashforth
scheme is used (see Figure~\ref{fig:particle_inertia_AB2}) to update the
particle position.
   \item The grid displacement in $y$-direction within a time step
    is defined by solving a Laplace equation
   \begin{equation}
      \nabla\cdot(k^e \nabla d_\text{g})= 0 \label{eq:laplace_ale4}
   \end{equation}
     where $d_\text{g}$ is equal to the particle $y$-displacement
     on the surface
     of the particle (curve \texttt{C5}). On all other boundaries 
     $d_\text{g}$ is set to zero.
     In Eq.~\eqref{eq:laplace_ale4} 
   $k^e$ is a function that controls the deformation of the domain.
   As suggested in \cite{Hu2001} we take $k^e$ 
   constant on an element and equal to the inverse of the
     area/volume of the element. In this way the deformation of the mesh is
largely performed by the big elements and the small ones around the solid
particle just move almost rigidly with the particle.
   \item The mesh velocity in $y$-direction $v^{n+1}_\text{g}$ at time 
    $t_{n+1}$ is defined by backwards
   differencing of the positions of the nodes of the mesh ($y_\text{g}$).
   The first time step we use a first-order differentiation:
\begin{equation}
  v^{n+1}_\text{g}=(y^{n+1}_\text{g} - y^{n}_\text{g})/\Delta t
\end{equation}
   and for steps two and further a second-order differentiation:
\begin{equation}
  v^{n+1}_\text{g}=(\frac32 y^{n+1}_\text{g} - 2 y^{n}_\text{g}
            +\frac12 y^{n-1}_\text{g})/\Delta t
\end{equation}
   \item The mesh velocity in $x$-direction $u^{n+1}_\text{g}$ at time 
    $t_{n+1}$ is constant in space and found by backwards
   differencing of the $x$-position of the particle.
   The first time step we use a first-order differentiation
   and for steps two and further a second-order (BDF2) differentiation.
   Due to this procedure the particle does not move in $x$-direction relative
   to the mesh.
  \item The inertia terms (of both fluid and particle)
        are time discretized using an implicit scheme (first
        step: Euler implicit, then second-order implicit Gear).
  \item The non-linear convection term for the inertia are solved using
        a Newton-Raphson iteration scheme in \texttt{ale5.f90} and
        \texttt{ale6.f90}. In \texttt{ale7.f90} and \texttt{ale8.f90}
         the inertia convection terms
        are discretized using a second-order predictor
        (Eq.~\ref{eq:predictor_inertia}), which avoids expensive iteration.
   \item 
     In \texttt{ale5.f90} the explicit stress formulation
     is used (see Sec.~\ref{sec:explicitstress})\footnote{Like the \texttt{ale1.f90} example, \texttt{ale5.f90} does not run currently.
Combining the ALE approach with time-integration schemes requiring
multiple evaluations of the convection term at different times in a single time
step does not take the proper mesh coordinates into account. Use examples \texttt{ale6.f90}, \texttt{ale7.f90} or \texttt{ale8.f90} instead.}. 
     In \texttt{ale6.f90}, \texttt{ale7.f90} and \texttt{ale8.f90}
     the implicit stress formulation
     is used (see Sec.~\ref{sec:implicitstress}). 
     The example \texttt{ale8.f90} uses, in addition, the projection of
     $\exp(\ten s)$ as given by Eq.~\eqref{eq:expsprojection}.
\end{itemize}

In Figure.~\ref{fig:Huang_compare} we compare the migration results with 
Huang et al.\ \cite{Huang97}.
\begin{figure}[htbp!]
   \begin{center}
   \includegraphics[scale=0.6]{Huang_compare}
   \end{center}
   \caption{Vertical position of the particle as a function of time. Comparison
    to the results in Huang et al.\ \cite{Huang97}. We find time and space
converged solutions for various methods, but there is still a difference with
\cite{Huang97}. Maybe the meshes in \cite{Huang97} are not refined enough.}
    \label{fig:Huang_compare}
\end{figure}


\subsection{Extrudate swell of a Newtonian fluid}
\label{sec:swellNewtonian}

(This example has been developed by Young Joon Choi of the TU/e.)

The example is derived from the examples in Sec.~\ref{sec:stick_slip}.
The slip surface is replaced by a free surface.
The fluid is Newtonian.
Use is made of the add-on surface advection (Sec.~\ref{sec:surfaceadvection}).
See \cite{Choi11c,Choi11b} and the references therein for the details on
solving the surface advection equation.

The example file is \texttt{extrudate\_swell1.f90} and can
be obtained by typing at the
command line:
\begin{quote}
   \texttt{getexample extrudate\_swell}
\end{quote}
The example uses an open boundary condition at the inflow boundary.
We assume zero normal traction and zero vertical velocity at the exit.
The free surface is traction free.
The flow is generated by imposing a flow rate constraint on the inflow
boundary.
The flow can be chosen to be planar of axisymmetric.
Below the free surface the nodes of the ALE mesh move proportional to the
vertical height of the free surface and the initial position of the
node.

You can run the example by
\begin{quote}
   \texttt{mesh\_two\_blocks < meshinput.txt}\\
   \texttt{extrudate\_swell1 < input\_extrudate\_swell1.txt}\\
   \texttt{post\_swell1}\\
   \texttt{streamfunction}
\end{quote}
Inspect the source files to see which data and plotfiles are generated.
Results of the final mesh and streamlines for a planar flow are shown in
Fig.~\ref{fig:swell_newtonian}.
\begin{figure}[htbp!]
   \begin{center}
   \includegraphics[scale=0.8]{mesh_swell_0200}\\[2mm]
   \includegraphics[scale=0.8]{streamfunction_0200}
   \end{center}
   \caption{Final mesh and streamlines
    of the extrudate swell problem using a Newtonian fluid.}
    \label{fig:swell_newtonian}
\end{figure}

In the example there is an option to add surface tension. The implementation is
based on a membrane with an isotropic surface stress $\gamma$. This means,
we add to the right-hand side of the weak form:
\begin{equation}
  -\int_S(\nabla\vek v)^T:\gamma(\ten I-\vek n\vek n)\,dS+
\int_{\partial S}\vek v\cdot\gamma\vek n_\text{c}\,ds
\end{equation}
where $\vek n$ is the unit normal vector on the surface $S$, $\partial S$ the
boundary of the surface (a point in 2D and a curve in 3D) and
$\vek n_\text{c}$ the
unit vector that is tangential to $S$ and normal to $\partial S$ and pointing
``outwards''. In the example \texttt{extrudate\_swell1.f90} the surface tension
integral is implemented using an object. A more simple approach is integration
on a curve (surface in 3D) and using $\nabla_\text{s}\vek v$ ($\nabla_\text{s}$
is the surface gradient operator) instead of $\nabla\vek v$. Furthermore, the
surface identity tensor ($\ten I-\vek n\vek n$) and $\nabla_\text{s}\vek v$ are
rewritten using the tangential (dual) base vectors. This approach is
used in \texttt{extrudate\_swell6.f90} and can be run by:
\begin{quote}
   \texttt{mesh\_two\_blocks < meshinput.txt}\\
   \texttt{extrudate\_swell6 < input\_extrudate\_swell1.txt}\\
   \texttt{post\_swell1}\\
   \texttt{streamfunction}
\end{quote}

\subsection{Extrudate swell of a viscoelastic fluid. Conventional ALE.}
\label{sec:swellviscoelastic}

(These examples have been developed by Young Joon Choi of the TU/e.)

The examples here are derived from the example \texttt{stick\_slip4.f90}
(second-order stress-implicit formulation) in Sec.~\ref{sec:stick_slip_impl}.
The slip surface is replaced by a free surface.
Use is made of the add-on surface advection (Sec.~\ref{sec:surfaceadvection}).
See \cite{Choi11c,Choi11b} and the references therein for the details on
solving the surface advection equation.

The example files are \texttt{extrudate\_swell2.f90} and
\texttt{extrudate\_swell3.f90} and can
be obtained by typing at the
command line:
\begin{quote}
   \texttt{getexample extrudate\_swell}
\end{quote}
The example uses an open boundary condition at the inflow boundary.
We assume zero normal traction and zero vertical velocity at the exit.
The free surface is traction free.
The flow is generated by imposing a flow rate constraint on the inflow
boundary.
Below the free surface the nodes of the ALE mesh move proportional to the
vertical height of the free surface and the initial position of the
node.
The time integration is second-order implicit Gear with conformation prediction
(Equation~\eqref{eq:Gearimp}).

The example \texttt{extrudate\_swell2.f90} uses a UCM/Oldroyd-B model in a
planar flow.
You can run the example by
\begin{quote}
   \texttt{mesh\_two\_blocks < meshinput.txt}\\
   \texttt{extrudate\_swell2 < input\_extrudate\_swell2.txt}\\
   \texttt{post\_swell2}\\
   \texttt{streamfunction}
\end{quote}
The example \texttt{extrudate\_swell3.f90} uses a Giesekus model in an
axisymmetric flow.
You can run the example by
\begin{quote}
   \texttt{mesh\_two\_blocks < meshinput.txt}\\
   \texttt{extrudate\_swell3 < input\_extrudate\_swell3.txt}\\
   \texttt{post\_swell3}\\
   \texttt{streamfunction}
\end{quote}
Inspect the source files to see which data and plotfiles are generated in the
examples.

\subsection{3D swell problem containing sharp edges}

(This example has been developed by Michelle Spanjaards of the TU/e.)

The examples \texttt{extrudate\_swell7.f90} and \texttt{extrudate\_swell8.f90} describe a transient 3D extrudate swell problem with a square die for a purely viscous fluid and a viscoelastic fluid, respectively. In many extrusion processes, the dies have complex shapes containing sharp edges such as squares. The sharp edge is maintained over a large distance in the extrudate and will swell in two directions. In this example the sharp edges are described as material lines. The positions of these lines can be used to describe the transverse swelling of the 2D free surfaces and expand the domain over which a 2D height function on the free surfaces is applied. Solving the 2D height functions gives the positions of the free surfaces \cite{spanjaards_transient_2019}. Figure \ref{fig:swellsquare} shows the swell in time for a quarter of an extrudate coming from a square die. Here, $t^{*} = H_{0}/U_{\text{avg}}$ with $H_{0}$ the initial half height of the die and $U_{\text{avg}}$ the average velocity over the inlet.

\begin{figure}[h!]
	
	\includegraphics[scale=.3]{squareN1.png}
	\includegraphics[scale=.3]{squareN2.png}
	\includegraphics[scale=.3]{squareN3.png}
	
	\caption{Swell of a quarter of a square extrudate in time (left to right): $t/t^{*}=0$, $6.25$ and $12.5$. }
	\label{fig:swellsquare}
\end{figure}

\subsection{Extrudate swell of a viscoelastic fluid. \xfem\ with temporary ALE.}
\label{sec:swellviscoelasticxfem}

(This example has been developed by Young Joon Choi of the TU/e.)

The example here is derived from the example \texttt{extrudate\_swell4.f90},
see Sec.~\ref{sec:swellNewtonian_xfem}
(Extrudate swell of a Newtonian fluid using \xfem).

The example file is \texttt{extrudate\_swell5.f90} and can
be obtained by typing at the
command line:
\begin{quote}
   \texttt{getexample extrudate\_swell}
\end{quote}
The example uses an open boundary condition at the inflow boundary.
We assume zero normal traction and zero vertical velocity at the exit.
The free surface is traction free.
The flow is generated by imposing a flow rate constraint on the inflow
boundary. The flow is planar.
The example uses a UCM/Oldroyd-B model.
The time integration is second-order implicit Gear with conformation prediction
(Equation~\eqref{eq:Gearimp}).
A temporary ALE scheme is used to handle the convection term in the
constitutive equation. For a description of the temporary ALE scheme, as
implemented in \tfem, see the papers \cite{Choi11c,choi_alignment_2012} and also the
thesis \cite{Choi11b}.
In the example there is an option to add surface tension (see
Sec.~\ref{sec:swellNewtonian_xfem}).

You can run the example by
\begin{quote}
   \texttt{mesh\_extrudate\_swell\_xfem < meshinput\_xfem.txt}\\
   \texttt{extrudate\_swell5 < input\_extrudate\_swell5.txt}\\
   \texttt{post\_swell5}
\end{quote}
Inspect the source files to see which data and plotfiles are generated.

\subsection{Spinning of a Newtonian fluid}
\label{sec:spinningNewtonian}

The example is derived from the extrudate swell
example in Sec.~\ref{sec:spinningNewtonian}.
The fluid is Newtonian.
Use is made of the add-on surface advection (Sec.~\ref{sec:surfaceadvection}).
See \cite{Choi11c,Choi11b} and the references therein for the details on
solving the surface advection equation.

The example file is \texttt{spinning1.f90} and can
be obtained by typing at the
command line:
\begin{quote}
   \texttt{getexample spinning}
\end{quote}
The example uses an open boundary condition at the inflow boundary.
We assume a prescribed
horizontal velocity and zero vertical velocity at the exit.
The free surface is traction free.
The flow is generated by imposing a flow rate constraint on the inflow
boundary.
The flow is axisymmetric.
Below the free surface the nodes of the ALE mesh move proportional to the
vertical height of the free surface and the initial position of the
node.

You can run the example by
\begin{quote}
   \texttt{mesh\_two\_blocks < meshinput.txt}\\
   \texttt{spinning1 < input\_spinning1.txt}\\
   \texttt{post\_spinning1}\\
   \texttt{streamfunction}
\end{quote}
Inspect the source files to see which data and plotfiles are generated.
Results of the spin velocity and streamlines for a planar flow are shown in
Fig.~\ref{fig:spinning_newtonian}.
\begin{figure}[htbp!]
   \begin{center}
   \includegraphics[scale=0.8]{velocity_color_spinning}\\[1cm]
   \includegraphics[scale=0.8]{streamfunction_spinning}
   \end{center}
   \caption{Spin velocity and streamlines
    of a spinning problem with draw ratio $\text{DR}=10$
    using a Newtonian fluid.}
    \label{fig:spinning_newtonian}
\end{figure}

\subsection{Spinning of a viscoelastic fluid}
\label{sec:spinningviscoelastic}

The examples here are derived from the extrudate swell examples
in Sec.~\ref{sec:swellviscoelastic}.
Use is made of the add-on surface advection (Sec.~\ref{sec:surfaceadvection}).
See \cite{Choi11c,Choi11b} and the references therein for the details on
solving the surface advection equation.

The example files are \texttt{spinning2.f90} and
\texttt{spinning3.f90} and can
be obtained by typing at the
command line:
\begin{quote}
   \texttt{getexample spinning}
\end{quote}
The example uses an open boundary condition at the inflow boundary.
We assume a prescribed
horizontal velocity and zero vertical velocity at the exit.
The free surface is traction free.
The flow is generated by imposing a flow rate constraint on the inflow
boundary.
Below the free surface the nodes of the ALE mesh move proportional to the
vertical height of the free surface and the initial position of the
node.
The time integration is second-order implicit Gear with conformation prediction
(Equation~\eqref{eq:Gearimp}).

The example \texttt{spinning2.f90} uses a Giesekus model.
You can run the example by
\begin{quote}
   \texttt{mesh\_two\_blocks < meshinput.txt}\\
   \texttt{spinning2 < input\_spinning2.txt}\\
   \texttt{post\_spinning2}\\
   \texttt{streamfunction}
\end{quote}
Results of the spin velocity and streamlines for $\text{DR}=10$ are shown in
Fig.~\ref{fig:spinning_viscoelastic}.
\begin{figure}[htbp!]
   \begin{center}
   \includegraphics[scale=0.8]{velocity_color_spinning2}\\[1cm]
   \hspace*{-1.1cm}\includegraphics[scale=0.8]{streamfunction_spinning2}
   \end{center}
   \caption{Spin velocity and streamlines
    of a spinning problem with draw ratio $\text{DR}=10$
    using a Giesekus fluid.}
    \label{fig:spinning_viscoelastic}
\end{figure}

The example \texttt{spinning3.f90} uses an XPP model.
You can run the example by
\begin{quote}
   \texttt{mesh\_two\_blocks < meshinput.txt}\\
   \texttt{spinning3 < input\_spinning3.txt}\\
   \texttt{post\_spinning2}\\
   \texttt{streamfunction}
\end{quote}
Inspect the source files to see which data and plotfiles are generated in the
examples. Note, that the example \texttt{spinning3.f90} uses the same post
program as \texttt{spinning2.f90} (for now).

\subsection{Oscillating planar Couette flow with a freely suspended particle
in a Newtonian liquid. Remeshing using gmsh.}
\label{sec:couette_ale_stokes_remeshing}

(These examples have been developed by Gaetano d'Avino of the University of
Naples)

We again
consider a planar Couette problem (shear flow between two moving walls) with
a (off-center) freely suspended particle. The example is similar to the example
in Sec.~\ref{sec:couette_ale}, but now the fluid is Newtonian (Stokes flow) and
the fluid domain is remeshed using gmsh (see Sec.~\ref{sec:readgmshmesh}),
either at regular intervals or if the mesh is too deformed.

The example files are \texttt{ale9.f90} (remeshing at regular time intervals)
and \texttt{ale10.f90} (remeshing if the mesh is too deformed)
and can be obtained by typing at the
command line:
\medskip
\begin{quote}
   \texttt{getexample ale}
\end{quote}

Some notes on the source:
\begin{itemize}
   \item The routine \texttt{generate\_read\_mesh} performs the remeshing. A general
  ``incomplete'' geometry file for gmsh (\texttt{particles\_in\_a\_box\_2D.igo}) is made available, which does not change.
  A geometry file (\texttt{mesh.geo}) is written with the actual data for the position of the particle,
  which `includes' the incomplete geometry file. The remeshing is performed by running gmsh
  from the command line using the intrinsic subroutine \texttt{execute\_command\_line}. 
  The new mesh is read from a file written by gmsh.
   \item Physical entities are used in gmsh to identify the boundary of the particle by a single curve.
   \item The grid displacement within a time step $\vek d_\text{g}$
    is defined by solving a Laplace equation
   \begin{equation}
      \nabla\cdot(k^e\nabla\vek d_\text{g})=\vek 0 \label{eq:laplace_ale3}
   \end{equation}
     $\vek d_\text{g}$ is equal to the particle displacement
     on the surface
     of the particle (curve \texttt{C5}). On all other boundaries $\vek
     d_\text{g}$ is set to zero. In Eq.~\eqref{eq:laplace_ale3} 
   $k^e$ is a function that controls the deformation of the domain.
   As suggested in \cite{Hu2001} we take $k^e$ 
   constant on an element and equal to the inverse of the
     area/volume of the element. In this way the deformation of the mesh is
largely performed by the big elements and the small ones around the solid
particle just move almost rigidly with the particle.
   \item The mesh deformation is measured by 
\[
   f^e_1 = |\log(A^e/A^e_0)|\quad\text{ and } \quad
   f^e_2 = |\log(S^e/S^e_0)|
\]
where $A^e$ and $S^e$ are the area and 
the aspect ratio of an element $e$, respectively.
The subindex 0 means the initial value.
The aspect ratio of an element $S^e$ is defined by
\[ 
   S^e=(L_\text{max}^e)^2/A^e
\]
where $L^e_\text{max}$ is the maximum length of the sides. 
Remeshing is performed if either of
\[
  f_1=\max_e f^e_1 \quad \text{ or } \quad 
  f_2=\max_e f^e_2 
\]
exceed a certain
threshold. This mesh quality criterion is identical to the one used
in \cite{Hu2001}, where a threshold of 1.39 is used. This corresponds to the
area or aspect ratio allowing to become four times larger or smaller than
initially before remeshing. In the example \texttt{ale10.f90} we use the same
threshold value of 1.39 (factor of four). In Figure~\ref{fig:remesh} we see the
initial mesh, the deformed mesh and the mesh after remeshing.
\end{itemize}
\begin{figure}[hb!]
   \begin{center}
   \includegraphics[scale=0.5]{mesh0000}\\[1cm]
   \includegraphics[scale=0.5]{mesh0034}\\[1cm]
   \includegraphics[scale=0.5]{mesh0035}
   \end{center}
   \caption{The initial mesh, a deformed mesh (after step 34)
            and a mesh after remeshing (after step 35).}
    \label{fig:remesh}
\end{figure}


\subsection{Couette flow with a freely suspended particle
in a viscoelastic liquid. Remeshing using gmsh.}
\label{sec:couette_ale_viscoelastic_remeshing}

(This example has been developed by Nick Jaensson of the TU/e).

We again
consider a Couette problem (shear flow between two moving walls) with
a (off-center) freely suspended particle. The example is similar to the example
in Sec.~\ref{sec:couette_ale_stokes_remeshing}, but now the fluid is viscoelastic.
The example files are \texttt{ale11.f90} and \texttt{ale12.f90} which use the
explicit stress formulation and implicit stress formulation respectively.
An extension of \texttt{ale12.f90} to 3D 
is available as example \texttt{ale13.f90}.

The time discretization used in \texttt{ale11.f90} (explicit stress) is
the semi-implicit Gear scheme with velocity prediction \cite{Karniadakis91,Hulsen96}.
\begin{equation}
\frac32 \frac{\ten c^{n+1}}{\Delta t}+(\hat{\vek u}^{n+1}-
 \vek u^{n+1}_\text{g})\cdot\nabla\ten c^{n+1}=
 \frac{2\ten c^{n}-\frac12\ten c^{n-1}}{\Delta t} +
2\ten g(\ten G^n,\ten c^n) -\ten g(\ten G^{n-1},\ten c^{n-1})
\end{equation}
where $\hat{\vek u}^{n+1}=2\vek u^n-\vek u^{n-1}$ and $\vek u^{n+1}_\text{g}$ is
the grid velocity. Since a displacement based mesh deformation approach is used, 
the mesh velocity can be determined using backwards differencing (BDF) and is thus 
known at the current time step $t_{n+1}$.
The velocity, pressure and velocity gradients at $t=0$ are found by
solving the velocity-pressure-gradient
problem without the extra stress tensor $\ten \tau$. The conformation tensor at $t=0$ is
initialized with zero stress. The time stepping scheme is started by using a
combined forward/backward Euler scheme with velocity prediction
\begin{equation}
 \frac{\ten c^{n+1}}{\Delta t}+(\hat{\vek u}^{n+1}-
 \vek u^{n+1}_\text{g}) \cdot\nabla\ten c^{n+1}=
 \frac{\ten c^{n}}{\Delta t} +\ten g(\ten G^n,\ten c^n)
\end{equation}
where $\hat{\vek u}^{n+1}=\vek u^n$. 

For the time discretization, some field variables need to be known at previous
timesteps. For the semi-implicit Gear scheme these are:
$\vek u^n$, $\vek u^{n-1}$, $\ten G^n$, $\ten G^{n-1}$, $\ten c^n$
and $\ten c^{n-1}$, after remeshing these will be obtained on the new mesh by solving a projection
problem. Note, that in the BDF schemes to determine the mesh velocity, the previous mesh
coordinates, $\vek x^n_\texttt{m}$ and $\vek x ^{n-1}_\texttt{m}$ are needed, which
are not known when a remeshing is performed. To find the mesh velocity after remeshing, the
new (remeshed) mesh at the previous time steps is found by projecting the coordinates of the
old (discarded) mesh at the previous time steps (see Fig.~\ref{fig:ale_remeshing}).

In \texttt{ale12.f90} (and \texttt{ale13.f90})
the implicit stress formulation is used, which means
that the velocity-pressure-gradient problem
is solved first in a time step, using a prediction for the conformation tensor.
This yields the velocity and velocity gradients
at the current time step, $\vek u^{n+1}$ and $\ten G^{n+1}$ which can be used
directly in the viscoelastic problem, so the velocities
and gradients do not have to be projected. 

Some notes on the source:
\begin{itemize}
   \item Similarly to \texttt{ale10.f90}, remeshing is performed when the
     mesh becomes too distorted. By default, the remeshing criterion is set unnecessarily low
     to show the effect of remeshing better.
   \item It would be possible to use the (known) Dirichlet boundary conditions
     and/or constraints in the
     projection problem (e.g. the velocity on the walls and the particle or periodic boundaries).
     Although this would increase the accuracy of the projection, the same system
     matrix can not be used for all projections. In \texttt{ale11.f90} and
     \texttt{ale12.f90}, no Dirichlet
     boundary conditions are used in the projection problem, and the same system matrix
     is used for all projections.
   \item The velocity and mesh coordinates are quadratically interpolated ($P_2$),
     so the projection is also quadratic. The conformation tensor and velocity gradients
     are linearly interpolated ($P_1$), so the projection is linear. To be able to project the 
     field variables using the same system matrix, a vector is created in each of the problems
     with the number of components to be projected. The variables \texttt{ncompv\_P2} and
     \texttt{ncompv\_P1} give the number of components to be projected using $P_2$ and $P_1$
     respectively. Be aware that the gradients are part of the 
     gradient-velocity-pressure problem, but are projected using a vector in the
     conformation problem. So when the data is transferred using \texttt{transfer\_data},
     both problems should be given.
   \item The logical \texttt{check\_projection} can be set to \texttt{.true.} to print VTK
    files before and after the projection that can be used to investigate the projection.
   \item For the integration in the projection problem, tabulated
    results using Gauss-Legendre are used. The integration rule can be set in 
    \texttt{gauss\_proj}
    which, for triangles, gives the maximum polynomial order for exact integration. 
\item Similarly to \texttt{ale8.f90}, if the log conformation approach is used (\texttt{logc = 1}), a
  projection problem is solved at each time step to find $\ten c^n$ from $\ten s^n$
  ($\ten c^n = \exp({\ten s^n})$).
\end{itemize}

\begin{figure}[htbp!]
   \begin{center}
   \includegraphics[scale=1.0]{ale_remeshing}
   \end{center}
   \caption{When a remeshing is performed, a new (remeshed)
     mesh (with coordinates $\vek x_\text{m,new}^{n+1}$)
     is created using the particle position $\vek X_\text{p}^{n+1}$. To obtain the new mesh at
     $t_n$ ($\vek x_\text{m,new}^n$), the coordinates of the old (discarded) mesh at
     $t_n$ ($\vek x_\text{m,old}^n$) are set as nodal values 
     in the old mesh at $t_{n+1}$ and projected onto the new mesh at $t_{n+1}$. The new mesh at 
     $t_n$ ($\vek x_\text{m,new}^n$)  can now simply be obtained by using the projected
     values as nodal coordinates. The new mesh at $t_{n-1}$ is obtained in a similar way. 
     Note, that the mesh deformation is exaggerated for clarity.}
    \label{fig:ale_remeshing}
\end{figure}

\subsection{Moving boundary due to capillary forces (surface tension, contact angle, gravity).}
\label{sec:capillary}

This example computes the interface between a Newtonian fluid in a cup
and gas (having a constant zero pressure). The surface tension is
$\gamma$, the density of the fluid is $\rho$ and the gravity constant is
$g$. Assuming a circular cup of radius $R$, the shape of the interface
(meniscus) is governed by the Bond number:
\[
   \mathrm{Bo} = \frac{\rho g R^2}\gamma
\]
For large values of the Bond number ($\mathrm{Bo}\gg1$) gravity dominates over surface tension and the interface is flat, except a region near the wall. The size of that wall region is given by the capillary length:
\[
    L_\text{c} = \sqrt{\frac{\gamma}{\rho g}}
\]
The height of the interface in this region is governed by the contact angle $\theta_\text{c}$, see Figure~\ref{fig:meniscus}. Note, that $\mathrm{Bo}=R^2/L_\text{c}^2$. The contact angle is given by Young's equation:
\[
  \gamma\cos\theta_\text{c}+\gamma_\text{fw}-\gamma_\text{aw}=0\quad\Rightarrow\quad 
  \cos\theta_\text{c}=-\frac{\gamma_\text{fw}-\gamma_\text{aw}}\gamma=-\frac{\Delta\gamma_\text{w}}\gamma
\]
where we introduced the difference in wall surface tension $\Delta\gamma_\text{w}=\gamma_\text{fw}-\gamma_\text{aw}.$
\begin{figure}[htbp!]
   \begin{center}
   \includegraphics[scale=0.5]{meniscus}
   \end{center}
   \caption{Fluid-air interface.}
    \label{fig:meniscus}
\end{figure}

For medium values of the Bond numbers ($\mathrm{Bo}=\mathcal{O}(1-10)$) the capillary length is similar to the length of the domain and the shape of the interface is affected in the whole region and capillary effects become important.

For small values of the Bond numbers ($\mathrm{Bo}\ll1$) surface tension dominates and the shape of the interface is approximately spherical. For the limit $\mathrm{Bo}=0$ the shape is spherical with a Laplace pressure of $2\Delta\gamma_\text{w}/R$ and a radius of curvature $-R\gamma/\Delta\gamma_\text{w}$.

The example file is \texttt{moving\_boundary1.f90} and can
be obtained by typing at the
command line:
\begin{quote}
   \texttt{getexample moving\_boundary}
\end{quote}
Use is made of the add-on surface advection (Sec.~\ref{sec:surfaceadvection}) to solve the height function equation.
The contact angle is imposed by applying a surface tension on the fluid-wall interface of $\Delta\gamma_\text{w}$.
The fluid is Newtonian and the viscosity $\mu$ determines the time for the surface to become steady.
Perfect slip is assumed at the wall(s).
Both 2D and axisymmetric cases can be solved. Both the symmetric and the full problem (two walls) can be solved. In the axisymmetric case, the full problem is for a fluid between two concentric cylinders.
Below the free surface the nodes of the ALE mesh move according to the solution of a Poisson equation. The nodes on the wall(s) move proportional to the vertical height of the free surface and the position of the node.

You can run the example by
\begin{quote}
   \texttt{mesh\_two\_blocks < meshinput1.txt}\\
   \texttt{moving\_boundary1 < input\_moving\_boundary1.txt}
\end{quote}
Inspect the source files to see which data and plot files are generated.
Results of steady interface for three values of the Bond number are shown in
Fig.~\ref{fig:capillary}.

\begin{figure}[htbp!]
   \begin{center}
   \includegraphics[scale=0.07]{cup1}
   \includegraphics[scale=0.07]{cup2}
   \includegraphics[scale=0.07]{cup3}
   \end{center}
   \caption{Steady shape of the interface for three values of the Bond number (left to right): $\mathrm{Bo}=160$, $16$ and $0$. Axisymmetric case where the symmetry line is the right boundary of the domain. Contact angle $\theta_\text{c}=60^\circ$.}
    \label{fig:capillary}
\end{figure}

\subsection{Moving boundary in a rotating cup due to centrifugal forces and gravity.}
\label{sec:centrifuge}

This example computes the interface between a Newtonian fluid in a cup and gas (having a constant zero pressure). 
The cup is rotating with a rotation rate $\omega$, 
 the density of the fluid is $\rho$ and the gravity constant is $g$. Assuming a circular cup of radius $R$, the shape of the interface without surface tension is found by a balance of static pressure due to gravity ($\rho g h$) and pressure due to the rotation ($\frac12\rho\omega^2r^2$):
\[
   h(r) = \frac{\omega^2 r^2}{2g}
\]

The example file is \texttt{moving\_boundary2.f90} and can
be obtained by typing at the
command line:
\begin{quote}
   \texttt{getexample moving\_boundary}
\end{quote}
Use is made of the add-on surface advection (Sec.~\ref{sec:surfaceadvection}) to solve the height function equation.
The fluid is Newtonian. Surface tension can be included, similar to the example in Sec.~\ref{sec:capillary}.
Perfect slip is assumed at the wall(s).
Both 2D and axisymmetric cases can be solved. Both the symmetric and the full problem (two walls) can be solved. In the axisymmetric case, the full problem is for a fluid between two concentric cylinders.
Below the free surface the nodes of the ALE mesh move according to the solution of a Poisson equation. The nodes on the wall(s) move proportional to the vertical height of the free surface and the position of the node.

You can run the example by
\begin{quote}
   \texttt{mesh\_two\_blocks < meshinput2.txt}\\
   \texttt{moving\_boundary2 < input\_moving\_boundary2.txt}
\end{quote}
Inspect the source files to see which data and plot files are generated.
Results of a steady interface is shown in
Fig.~\ref{fig:centrifuge}. Note, that if there is not enough fluid in the cup to create the required height, the problem becomes much more difficult due to the thin fluid layer at the bottom of the cup.

\begin{figure}[htbp!]
   \begin{center}
   \includegraphics[scale=0.2]{centrifuge1}
   \includegraphics[scale=0.2]{centrifuge2}
   \end{center}
   \caption{Steady shape of the interface for $\omega=3$, $g=10$, $R=4$ (figure on the right). 
   The figure on the left is the initial fluid shape. Axisymmetric case where the symmetry line is the right boundary of the domain. Surface tension has been included (contact angle $\theta_\text{c}=60^\circ$).}
    \label{fig:centrifuge}
\end{figure}

\subsection{Moving boundary in a Couette device due to viscoelastic normal stress differences (Weissenberg effect) and gravity.}
\label{sec:weisseff}

This example computes the interface between a viscoelastic fluid in a Couette device and gas (having a constant zero pressure). 
The inner cylinder is rotating with a rotation rate $\omega$, 
 the density of the fluid is $\rho$ and the gravity constant is $g$.  
The example file is \texttt{moving\_boundary3.f90} and can
be obtained by typing at the
command line:
\begin{quote}
   \texttt{getexample moving\_boundary}
\end{quote}
Use is made of the add-on surface advection (Sec.~\ref{sec:surfaceadvection}) to solve the height function equation.
The fluid is viscoelastic (Giesekus). Surface tension can be included, similar to the example in Sec.~\ref{sec:capillary}. At the ``bottom'' there is perfect slip for radial and tangential motion. The vertical velocity at the bottom is set to zero (Dirichlet). At the inner and outer cylinder perfect slip is assumed in vertical direction. In radial and tangential direction the velocities are imposed (Dirichlet).
Below the free surface the nodes of the ALE mesh move according to the solution of a Poisson equation. The nodes on the wall(s) move proportional to the vertical height of the free surface and the position of the node.

You can run the example by
\begin{quote}
   \texttt{mesh\_two\_blocks < meshinput3.txt}\\
   \texttt{moving\_boundary3 < input\_moving\_boundary3.txt}
\end{quote}
Inspect the source files to see which data and plot files are generated.
Results of a steady interface is shown in
Fig.~\ref{fig:weisseffect}. 
A small amount of surface tension has been included to stabilize the interface. During the start-up of the flow
the interface seems to become vertical at some point, which cannot be handled by the height function approach.
 For this case, we have added a surface tension of $\gamma=2\Delta x G$, where $\Delta x$ is the size of elements on the interface and the elastic modulus is $G=\eta_\text{p}/\lambda$. The argument here is: we want the surface tension to be active on grid scale, i.e.\ the grid Laplace pressure $\gamma/2\Delta x\approx G$. The factor 2 seems to work for this example, but for other examples a different factor might be necessary. Note, that the grid emulsion time scale $\Delta x \eta_\text{p}/\gamma=\lambda/2$.
\begin{figure}[htbp!]
   \begin{center}
   \includegraphics[scale=0.15]{weisseff1}
   \includegraphics[scale=0.15]{weisseff2}
   \end{center}
   \caption{Steady shape of the interface for $\omega=1$, $\rho g=1$, $R_1=1$, $R_2=4$ (figure on the right). 
   Giesekus model $\eta_\text{s}=0.1$, $\eta_\text{p}=1$, $\lambda=2$, $\alpha=0.01$.
   The figure on the left is the initial fluid shape. Axisymmetric, where the inner cylinder (radius $R_1$) is at the right boundary of the domain. Surface tension has been included (contact angle $\theta_\text{c}=90^\circ$).}
\label{fig:weisseffect}
\end{figure}

There are two other examples, \texttt{moving\_boundary3a.f90} and \texttt{moving\_boundary3b.f90}, that use a different discretization of the velocities and/or height function. The example \texttt{moving\_boundary3a.f90}
uses $P_1$ shape functions for the height function, whereas \\
\texttt{moving\_boundary3b.f90} optionally uses $Q_1\text{-iso-}Q_2$ shape functions for the velocity and/or 
$P_1\text{-iso-}P_2$ line elements for the height function. Note however, that for this particular problem there is no need to use a different discretization and the 
original \texttt{moving\_boundary3.f90} is just fine.

\subsection{Open channel flow: curved boundary due to the second normal stress difference.}
\label{sec:curvedN2}

This example computes the viscoelastic flow in an open channell. The surface becomes curved due to second normal stress differences \cite{kuo_use_1974}. 
The example file is \texttt{moving\_boundary4.f90} and can
be obtained by typing at the
command line:
\begin{quote}
   \texttt{getexample moving\_boundary}
\end{quote}
Use is made of the add-on surface advection (Sec.~\ref{sec:surfaceadvection}) to solve the height function equation.
The fluid is viscoelastic (Giesekus). There is perfect slip on the side walls in vertical direction. The flow is imposed by a flow rate in third direction.

You can run the example by
\begin{quote}
   \texttt{mesh\_two\_blocks < meshinput4.txt}\\
   \texttt{moving\_boundary4 < input\_moving\_boundary4.txt}
\end{quote}
Inspect the source files to see which data and plot files are generated.
Results of a steady interface is shown in
Fig.~\ref{fig:curvedN2}. 
A small amount of surface tension has been included to stabilize the interface near the walls. See Sec.~\ref{sec:weisseff} for a discussion on the amount of surface tension added.
\begin{figure}[htbp!]
   \begin{center}
   \includegraphics[scale=0.1]{gootje1}
   \includegraphics[scale=0.1]{gootje2}
   \includegraphics[scale=0.1]{gootje3}
   \end{center}
   \caption{Steady shape of the interface for $\alpha=0$ (Oldroyd-B), $\alpha=0.01$ and $\alpha=0.1$ for the
   Giesekus model.
   Symmetry is assumed at the right boundary of the domain. Surface tension has been included for stability (contact angle $\theta_\text{c}=90^\circ$).}
\label{fig:curvedN2}
\end{figure}

\subsection{Fountain flow of a Newtonian fluid. Imposed parabolic in/out profile.}
\label{sec:fountainN1}

This example computes the fountain flow for a Newtonian fluid with an imposed in/out profile (see Fig.~\ref{fig:fountaingeom1}). See \cite{bogaerds_stability_2004} for more info on the fountain flow problem.
\begin{figure}[htbp!]
   \begin{center}
   \includegraphics[scale=1.0]{fountaingeom}
   \end{center}
   \caption{Flow domain for the fountain flow computation. The upper and lower wall have a velocity $U$ to the left. On the left boundary a parabolic velocity profile is imposed have a flow rate $Q=0$. On the right, the initally flat free surface will become curved and the contact points on the upper and lower walls will move to the left until a steady state is reached. Note: $H$ is the half-height of the domain and $L$ is the initial length of the domain with a flat free surface.}
\label{fig:fountaingeom1}
\end{figure}

The example file is \texttt{fountain\_flow1.f90} and can
be obtained by typing at the
command line:
\begin{quote}
   \texttt{getexample fountain\_flow}
\end{quote}
Use is made of the add-on surface advection (Sec.~\ref{sec:surfaceadvection}) to solve the height function equation for obtaining the position of the free surface. Taylor-Hood $Q_2/Q_1$ quad elements are used for the velocity/pressure and $P_2$ shape functions are used for the height function.

You can run the example by
\begin{quote}
   \texttt{mesh\_four\_blocks < meshinput1.txt}\\
   \texttt{fountain\_flow1 < input\_fountain\_flow1.txt}
\end{quote}
Inspect the source files to see which data and plot files are generated.

After starting the flow, the contact points will move to the left until a steady state is reached. In Figure~\ref{fig:contactpoint} the position of the contact point (upper and lower point are identical within truncation errors) is shown relative to the steady state value. We clearly see the a single exponential decay to a steady value over a large time interval.
\begin{figure}[htbp!]
   \begin{center}
   \includegraphics[scale=1.0]{contact_position}
   \end{center}
   \caption{Position of the contact point relative to the steady state value as a function of time. Note: $H=1$, $U=1$, $\mu=1$, $L=10$. The green line represents $\exp(-1.2t)$, where $t$ is the time.}
\label{fig:contactpoint}
\end{figure}
Results of a steady interface with streamlines is shown in
Fig.~\ref{fig:flowfront}. 
\begin{figure}[htbp!]
   \begin{center}
   \includegraphics[scale=0.2]{flowfront1}
   \end{center}
   \caption{Steady fountain flow with streamlines near the flow front.}
\label{fig:flowfront}
\end{figure}

There are two other examples, \texttt{fountain\_flow1a.f90} and \texttt{fountain\_flow1b.f90}, that use a different discretization of the velocities and/or height function. The example \\
\texttt{fountain\_flow1a.f90}
uses $P_1$ shape functions for the height function, whereas \\
\texttt{fountain\_flow1b.f90} uses $Q_1\text{-iso-}Q_2$ shape functions for the velocity and 
$P_1\text{-iso-}P_2$ line elements for the height function. Note however, that for this particular problem there is no need to use a different discretization, but it can be useful for viscoelastic fluid flow.

\subsection{Cone-plate with free surface (axisymmetric with swirl)}
\label{sec:cone-plate_ve_si_fs}

(This example has been contributed by Michelle Spanjaards of the TU/e)

The problem is an extension of the viscoelastic problem in Sec.~\ref{sec:cone-plate_ve_si}.
The geometry is depicted in Figure~\ref{fig:coneplategeom}. The plate (curve 1) is stationary ($\vek u=\vek 0$). The cone (curve 2) is rotating around the $z$-axis ($\vek u=\omega r \vek e_\theta$). Curve 3 is a traction free moving surface ($\vek t_\text{N}=\vek 0$) and is \emph{initially} of a circular shape having point 1 as its center. The angle between the cone and plate $\theta_\text{cp}$ is 0.1~rad in the figure, but can be adjusted. The shear rate is approximately constant and is given by $\omega/\theta_\text{cp}$. 

The initially circular shape of the free surface cannot be maintained and the surface starts to move.
The contact points 2 and 3 are able to move along the plate and cone, respectively. It is assumed that there is perfect slip (no shear stress) in a direction tangential to the plane and normal the $\theta$ direction of Figure~\ref{fig:coneplategeom}. Optionally, surface tension can be added.

The free surface is described in a polar coordinate system\footnote{Actually, the free surface is described in a \emph{spherical} coordinate system $(\phi,\rho,\theta)$, where the polar system is obtained for the cross section at a fixed $\theta$.} $(\phi,\rho)$, where $\phi$ is the angle relative to the plate and $\rho$ the distance to the origin, see Figure~\ref{fig:polar_cone_plate}. The height function describing the free surface is taken to be the distance from the origin as a function of the angle $\phi$,
i.e.\ $\rho(\phi)$, for the points on the interface. The kinematic equation for $\rho(\phi)$ to be solved on a fixed domain $\phi\in(0,\theta_\text{cp})$ is given by
\[
   \pderiv{\rho}t+\frac{u_\phi}{\rho}\pderiv{\rho}\phi=u_\rho
\]
where $u_\phi=\vek e_{\phi}\cdot\vek u$ and $u_\rho= \vek e_\rho\cdot \vek u$.

\begin{figure}[htbp!]
   \begin{center}
   \includegraphics[scale=0.8]{polar}
   \end{center}
   \caption{Polar coordinate system used for the height function.}
\label{fig:polar_cone_plate}
\end{figure}
The source of the example is in the file \texttt{coneplate4.f90}
and can be obtained by typing at the
command line:
\begin{quote}
   \texttt{getexample coneplate}
\end{quote}

\subsection{Sagging of a cylindrical fluid column between two flat plates}
\label{sec:sagging}

This example computes the sagging of cylindrical column between two flat plates under the influence of gravity. 
Surface tension at the fluid-air interface and a contact angle between the fluid and wall has been taken into account. Perfect slip (zero friction) has been assumed between the fluid and flat plates.

The example file is \texttt{moving\_boundary5.f90} and can
be obtained by typing at the
command line:
\begin{quote}
   \texttt{getexample moving\_boundary}
\end{quote}
Use is made of the add-on surface advection (Sec.~\ref{sec:surfaceadvection}) to solve the height function equation.
The fluid is Newtonian and inertia has been neglected.

You can run the example by
\begin{quote}
   \texttt{mesh\_two\_blocks < meshinput5.txt}\\
   \texttt{moving\_boundary5 < input\_moving\_boundary5.txt}
\end{quote}
Inspect the source files to see which data and plot files are generated.
Results are shown in Fig.~\ref{fig:sagging}. Note, that the initial column is not consistent with the imposed contact angle and that there is an immediate shape change at $t=0^+$.
\begin{figure}[htbp!]
   \begin{center}
   \includegraphics[scale=0.08]{sagging1}
   \includegraphics[scale=0.08]{sagging2}
   \end{center}
   \caption{Sagging of an initially cylindrical fluid column under the influence of gravity. Left: Initial column. Right: Column after some time.}
\label{fig:sagging}
\end{figure}

